
\section{AT-01: Finite nuclear size in a hydrogen-like ion}
\source{2026 Atomic Assignment 3 (\texttt{PHYS4055\_A3.pdf}), Question 1(a).}
\subsection*{Step 1: construct the perturbing potential}
In atomic units (\(a_0=1,\ e=\hbar=m_e=1\)),
a point nucleus has potential \(V_0(r)=-Z/r\).
For uniform nuclear charge in a sphere of radius \(R\),
\[
V_R(r)=
\begin{cases}
-\dfrac{Z}{2R}\left(3-\dfrac{r^2}{R^2}\right),&r<R,\\[1mm]
-\dfrac Zr,&r\ge R .
\end{cases}
\]
The perturbation is therefore
\[
\Delta V(r)=
\begin{cases}
\displaystyle Z\!\left(\frac1r-\frac{3}{2R}
+\frac{r^2}{2R^3}\right),&r<R,\\
0,&r\ge R .
\end{cases}
\]
It is nonnegative inside the nucleus: a distributed
charge is less singular than a point charge.

\subsection*{Step 2: use first-order perturbation theory}
For sufficiently small \(R\), approximate an \(ns\)
wavefunction as constant inside the nucleus.
Its nucleus density for a hydrogenic state is
\(|\psi_{ns}(0)|^2=Z^3/(\pi n^3)\).
Then
\[
\Delta E_{ns}^{(1)}
\simeq4\pi|\psi_{ns}(0)|^2
\int_0^R r^2\Delta V(r)\,dr .
\]
Do the integral explicitly:
\begin{align*}
\int_0^R r^2\Delta V(r)\,dr
&=Z\int_0^R
\left(r-\frac{3r^2}{2R}+\frac{r^4}{2R^3}\right)dr\\
&=ZR^2\left(\frac12-\frac12+\frac1{10}\right)
=\frac{ZR^2}{10}.
\end{align*}
Thus
\[
\boxed{\Delta E_{\rm fns}(ns)\simeq
\frac{2Z^4R^2}{5n^3}
=\frac{2Z^2}{5n^3}(ZR)^2\quad\text{(Hartree)}}.
\]
The result is positive and scales as \(R^2\);
it becomes negligible for a point nucleus.
The radius \(R\) here is the \emph{uniform-sphere} radius,
\(R=\sqrt{5/3}\,R_{\rm rms}\).

\subsection*{Physical check and limitation}
For \(l>0\), the wavefunction vanishes at the origin in
the nonrelativistic central-field model, and the leading
small-\(R\) correction is suppressed relative to \(s\)-states.
For high-\(Z\) ions, relativistic wavefunctions and more
accurate nuclear charge models may be needed.
\checkmethod{Write the \emph{difference} between finite-
and point-nucleus potentials, not the full Coulomb energy.
Use consistent atomic units and the correct radial measure \(r^2dr\).}
